% !TEX root = ../algebra_02.tex

\section{Поле разложения и конечные поля}

\subsection{Поле разложения}

\begin{Theorem}{Существование поля разложения}{}
	Пусть $f \in K[x]$. Тогда существует поле разложения многочлена $f$.
\end{Theorem}

\begin{proof}
	Докажем индукцией по $\deg f$.
	
	\textbf{База:} При $\deg f = 1$ в качестве поля разложения подходит само поле $L = K$.
	
	\textbf{Шаг:} Пусть $\deg f > 1$. Пусть $f_0 \mid f$ --- неприводимый делитель в кольце $K[x]$. Построим факторкольцо (которое является полем в силу неприводимости $f_0$):
	\[ K_1 := K[x]/(f_0) \]
	Обозначим класс $x$ в $K_1$ как $x_1 = \overline{x}$. Тогда по определению $f_0(x_1) = 0$, откуда следует, что и $f(x_1) = 0$. Значит, по теореме Безу в $K_1[x]$ многочлен $f$ представим в виде:
	\[ f = (x - x_1)g, \quad g \in K_1[x] \]
	По предположению индукции, примененному к многочлену $g \in K_1[x]$, существует поле разложения $L/K_1$. Тогда в $L[x]$ многочлен $g$ раскладывается на линейные множители:
	\[ g = \varepsilon(x - x_2) \dots (x - x_n) \]
	Следовательно, в $L[x]$ многочлен $f$ также полностью раскладывается:
	\[ f = \varepsilon(x - x_1)(x - x_2) \dots (x - x_n) \]
	Само поле разложения порождается всеми корнями:
	\[ L = K_1(x_2, \dots, x_n) = K(x_1)(x_2, \dots, x_n) = K(x_1, x_2, \dots, x_n) \]
\end{proof}

\subsection{Конечные поля}

Основные вопросы теории конечных полей: существование ($\exists$), единственность ($!$), описание автоморфизмов и подполей.

Пусть $F$ --- конечное поле. Тогда его характеристика отлична от нуля: $\operatorname{char} F = p > 0$, где $p$ --- простое число. Поле $F$ можно рассматривать как расширение простого подполя $\mathbb{F}_p = \mathbb{Z}/p\mathbb{Z}$. 
Пусть степень этого расширения равна $[F : \mathbb{F}_p] = n$. Тогда как линейное пространство над $\mathbb{F}_p$ поле $F$ изоморфно $\mathbb{F}_p^n$. Следовательно, количество элементов в конечном поле равно $|F| = p^n$.

\begin{Proposition}{Эндоморфизм Фробениуса}{}
	Пусть $F$ --- произвольное поле характеристики $p > 0$ (не обязательно конечное). Тогда отображение $\Phi: F \to F$, заданное правилом $x \mapsto x^p$, является эндоморфизмом поля $F$.
\end{Proposition}

\begin{proof}
	Проверим свойства гомоморфизма:
	\begin{enumerate}
		\item $\Phi(1) = 1^p = 1$.
		\item $\Phi(xy) = (xy)^p = x^p y^p = \Phi(x)\Phi(y)$.
		\item Для суммы по формуле бинома Ньютона имеем:
		\[ \Phi(x+y) = (x+y)^p = x^p + y^p + \sum_{i=1}^{p-1} C_p^i x^i y^{p-i} \]
		Так как $p$ --- простое, все биномиальные коэффициенты $C_p^i$ при $1 \le i \le p-1$ делятся на $p$, а значит, в поле характеристики $p$ они равны нулю. Получаем:
		\[ \Phi(x+y) = x^p + y^p = \Phi(x) + \Phi(y) \]
	\end{enumerate}
\end{proof}

\begin{Remark}{Инъективность гомоморфизмов полей}{}
	Любой гомоморфизм полей $\alpha: K \to L$ инъективен. Действительно, $\ker \alpha$ является идеалом в поле $K$, следовательно, $\ker \alpha = \{0\}$ или $\ker \alpha = K$. Так как $\alpha(1) = 1 \neq 0$, ядро не может совпадать со всем полем, поэтому $\ker \alpha = \{0\}$.
\end{Remark}

\begin{Corollary}{}{}
	Если $F$ --- конечное поле, то эндоморфизм Фробениуса $\Phi$ является автоморфизмом (так как инъективное отображение конечного множества в себя биективно).
\end{Corollary}

\subsection{Существование конечного поля и неприводимых многочленов}

Для построения поля из $p^n$ элементов как факторкольца $\mathbb{F}_p[x]/(f_0)$ требуется доказать существование неприводимого многочлена $f_0$ степени $n$. Докажем это через построение поля разложения.

Пусть $K$ --- поле разложения многочлена $X^{p^n} - X$ над полем $\mathbb{F}_p$. 
Определим множество $F$ как множество всех корней этого многочлена в $K$:
\[ F := \{ x \in K \mid x^{p^n} - x = 0 \} \]

Вычислим формальную производную многочлена $f = X^{p^n} - X$:
\[ f' = p^n X^{p^n - 1} - 1 = -1 \]
Так как производная не имеет общих корней с $f$, многочлен $f$ не имеет кратных корней. Следовательно, количество корней в точности равно его степени: $|F| = p^n$.

Покажем, что $F$ является полем. Пусть $x, y \in F$:
\begin{enumerate}
	\item $(x+y)^{p^n} = \Phi^n(x+y) = \Phi^n(x) + \Phi^n(y) = x^{p^n} + y^{p^n} = x + y$, откуда $x+y \in F$.
	\item $(xy)^{p^n} = x^{p^n}y^{p^n} = xy$, откуда $xy \in F$.
	\item Очевидно, что $1 \in F$.
\end{enumerate}
Таким образом, $F$ --- подполе в $K$, состоящее из $p^n$ элементов. Так как $K$ порождается корнями $f$, то $K = \mathbb{F}_p(F) = F$. Существование поля из $p^n$ элементов доказано.

\begin{Proposition}{О примитивном элементе}{}
	Пусть $|F| = p^n$. Тогда существует элемент $a \in F$ такой, что $F = \mathbb{F}_p(a)$.
\end{Proposition}

\begin{proof}
	Мультипликативная группа конечного поля $F^*$ конечна ($|F^*| = p^n - 1$). Известно, что любая конечная подгруппа мультипликативной группы поля является циклической. 
	Следовательно, $F^* = \langle a \rangle$ для некоторого элемента $a \in F^*$. 
	Тогда подполе $\mathbb{F}_p(a)$ содержит все степени $a$, то есть содержит всё множество $F^*$, а значит, $\mathbb{F}_p(a) = F$.
\end{proof}

\begin{Corollary}{Существование неприводимого многочлена любой степени}{}
	Для любого простого $p$ и любого $n \in \mathbb{N}$ существует неприводимый многочлен $f \in \mathbb{F}_p[x]$ степени $n$.
\end{Corollary}

\begin{proof}
	По доказанному выше, существует поле $\mathbb{F}_{p^n}$, и его можно представить как $\mathbb{F}_p(a)$ для некоторого $a$. Тогда минимальный многочлен элемента $a$ над $\mathbb{F}_p$ является неприводимым и имеет степень $n$.
\end{proof}

\begin{Remark}{}{}
	Если $f \in \mathbb{F}_p[x]$ --- неприводимый многочлен степени $n$, то $f \mid (X^{p^n} - X)$.
\end{Remark}

\begin{proof}
	Пусть $K = \mathbb{F}_p[x]/(f)$. Так как $\deg f = n$, то $|K| = p^n$. 
	Порядок мультипликативной группы равен $|K^*| = p^n - 1$. По следствию из теоремы Лагранжа, для любого $a \in K^*$ выполнено $a^{p^n - 1} = 1$. 
	Умножая на $a$, получаем, что для любого $a \in K$ выполнено $a^{p^n} = a$.
	В частности, для элемента $x = \overline{X} \in K$ выполнено $x^{p^n} - x = 0$. Это равносильно тому, что многочлен $X^{p^n} - X$ делится на $f$ в кольце $\mathbb{F}_p[x]$.
\end{proof}

\begin{Example}{Разложение $X^{p^n} - X$ для $p=2, n=2$}{}
	Рассмотрим многочлен $X^4 - X$ над полем $\mathbb{F}_2$:
	\[ X^4 - X = X(X^3 - 1) = X(X - 1)(X^2 + X + 1) \]
	Здесь $X^2 + X + 1$ --- единственный неприводимый многочлен степени 2 над $\mathbb{F}_2$.
\end{Example}

\subsection{Единственность конечного поля}

\begin{Theorem}{Единственность конечного поля с точностью до изоморфизма}{}
	Пусть $K$ и $L$ --- конечные поля из $p^n$ элементов. Тогда $K \cong L$.
\end{Theorem}

\begin{proof}
	По доказанному ранее, поле $K$ можно представить в виде $K = \mathbb{F}_p[x]/(f)$, где $f \in \mathbb{F}_p[x]$ --- неприводимый многочлен степени $n$. По замечанию о неприводимых многочленах, $f \mid (X^{p^n} - X)$. 
	
	В кольце многочленов $L[x]$ многочлен $X^{p^n} - X$ имеет ровно $p^n$ корней (поскольку для любого ненулевого элемента $a \in L^*$ выполнено $a^{p^n-1} = 1$, а ноль также является корнем). Следовательно, $X^{p^n} - X$ полностью раскладывается в $L[x]$ на линейные множители. Из делимости следует, что $f$ также раскладывается на линейные множители над $L$.
	
	Пусть $b \in L$ --- некоторый корень многочлена $f$. Рассмотрим простое расширение $\mathbb{F}_p(b) / \mathbb{F}_p$. Минимальным многочленом элемента $b$ над $\mathbb{F}_p$ является $f$, поэтому $[\mathbb{F}_p(b) : \mathbb{F}_p] = \deg f = n$. 
	Следовательно, $|\mathbb{F}_p(b)| = p^n$. Так как $\mathbb{F}_p(b) \subset L$ и оба поля состоят из $p^n$ элементов, то $\mathbb{F}_p(b) = L$.
	
	Итоговый изоморфизм строится следующим образом:
	\[ L = \mathbb{F}_p(b) \cong \mathbb{F}_p[x]/(f) = K \]
\end{proof}

\begin{Remark}{Обозначение конечного поля}{}
	Поскольку конечные поля одного порядка изоморфны, поле из $p^n$ элементов обозначается стандартно как $\mathbb{F}_{p^n}$.
\end{Remark}

\begin{Example}{Изоморфизм полей из 8 элементов}{}
	Рассмотрим два представления поля из 8 элементов: $\mathbb{F}_2[x]/(x^3+x+1)$ и $\mathbb{F}_2[y]/(y^3+y^2+1)$. Изоморфизм между ними задается отображением $x \mapsto y+1$, поскольку $y+1$ является корнем многочлена $X^3+X+1$ во втором поле.
\end{Example}

\subsection{Автоморфизмы и подполя конечных полей}

\begin{Theorem}{Группа автоморфизмов конечного поля}{}
	Группа автоморфизмов $\operatorname{Aut}(\mathbb{F}_{p^n})$ является циклической группой порядка $n$, порожденной автоморфизмом Фробениуса $\Phi$.
\end{Theorem}

\begin{proof}
	Докажем в два этапа.
	\begin{enumerate}
		\item Покажем, что $\operatorname{ord} \Phi = n$. \\
		Для любого $a \in \mathbb{F}_{p^n}$ имеем $\Phi^n(a) = a^{p^n} = a$, откуда $\Phi^n = \operatorname{id}_{\mathbb{F}_{p^n}}$. 
		Допустим, существует $m < n$ такое, что $\Phi^m = \operatorname{id}_{\mathbb{F}_{p^n}}$. Тогда для любого $a \in \mathbb{F}_{p^n}$ выполнено $a^{p^m} - a = 0$. Получается, что многочлен $X^{p^m} - X$ имеет $p^n$ различных корней. Это невозможно, так как его степень $p^m$ строго меньше $p^n$. Значит, порядок $\Phi$ в точности равен $n$.
		
		\item Докажем, что $|\operatorname{Aut}(\mathbb{F}_{p^n})| \le n$. \\
		Предположим противное: существуют различные автоморфизмы $\psi_1, \dots, \psi_{n+1} \in \operatorname{Aut}(\mathbb{F}_{p^n})$.
		Мультипликативная группа конечного поля циклическая, поэтому $\mathbb{F}_{p^n}^* = \langle a \rangle$ для некоторого элемента $a$. Пусть $f = \operatorname{Irr}_{\mathbb{F}_p}(a)$ --- его минимальный многочлен над простым подполем $\mathbb{F}_p$, причем $\deg f = n$.
		
		Заметим, что любой автоморфизм $\psi_j$ сохраняет элементы $\mathbb{F}_p$ неподвижными (так как $\psi_j(1)=1$). Поскольку $f(a) = 0$ и коэффициенты $f$ лежат в $\mathbb{F}_p$, то:
		\[ \forall j : \quad \psi_j(f(a)) = 0 \implies f(\psi_j(a)) = 0 \]
		Таким образом, элементы $\psi_1(a), \dots, \psi_{n+1}(a)$ являются корнями многочлена $f$. Так как степень $f$ равна $n$, по принципу Дирихле найдутся индексы $i \neq j$, для которых $\psi_i(a) = \psi_j(a)$. 
		Но элемент $a$ порождает все поле $\mathbb{F}_{p^n}$, поэтому совпадение автоморфизмов на порождающем элементе влечет их полное совпадение: $\psi_i = \psi_j$. Мы получили противоречие.
	\end{enumerate}
\end{proof}

\begin{Theorem}{Структура подполей конечного поля}{}
	Пусть $p$ --- простое число, $m, n \in \mathbb{N}$.
	\begin{enumerate}
		\item Если $m \nmid n$, то в поле $\mathbb{F}_{p^n}$ нет подполей из $p^m$ элементов.
		\item Если $m \mid n$, то в поле $\mathbb{F}_{p^n}$ существует единственное подполе из $p^m$ элементов.
	\end{enumerate}
\end{Theorem}

\begin{proof}
	\begin{enumerate}
		\item Пусть $K \subset \mathbb{F}_{p^n}$ --- подполе, $|K| = p^m$. Тогда $\mathbb{F}_{p^n}$ можно рассматривать как линейное пространство над полем $K$. Пусть его размерность равна $[\mathbb{F}_{p^n} : K] = d \in \mathbb{N}$. Тогда $|\mathbb{F}_{p^n}| = |K|^d$, то есть $p^n = (p^m)^d = p^{md}$. Следовательно, $n = md$, откуда $m \mid n$. Противоречие.
		
		\item Докажем единственность. Пусть $K \subset \mathbb{F}_{p^n}$ --- подполе из $p^m$ элементов. Тогда для любого $a \in K$ выполнено $a^{p^m} - a = 0$. Следовательно, $K$ полностью содержится во множестве корней многочлена $X^{p^m}-X$:
		\[ K \subset \{ a \in \mathbb{F}_{p^n} \mid a^{p^m} - a = 0 \} = F \]
		Поскольку многочлен степени $p^m$ имеет не более $p^m$ корней, а множество $K$ содержит ровно $p^m$ элементов, то $K = F$. Значит, если подполе существует, оно состоит ровно из корней этого многочлена и потому единственно. Само множество $F$, как было доказано ранее, образует поле, что доказывает и существование.
	\end{enumerate}
\end{proof}

\subsection{Соответствие Галуа}

\begin{Definition}{Группа Галуа}{def_galois}
	Пусть $L/K$ --- конечное расширение полей. Группой автоморфизмов $L$ над $K$ называется множество всех автоморфизмов поля $L$, оставляющих неподвижными элементы подполя $K$:
	\[ \operatorname{Aut}(L/K) = \{ \sigma \in \operatorname{Aut}(L) \mid \sigma|_K = \operatorname{id}_K \} \]
\end{Definition}

\begin{Statement}{Соответствие Галуа для конечных полей}{}
	Существует взаимно однозначное соответствие (две взаимно обратные биекции) между подполями $F$ расширения (где $K \subset F \subset L$) и подгруппами $H \le \operatorname{Aut}(L/K)$:
	\begin{itemize}
		\item Каждому подполю $F$ ставится в соответствие подгруппа $H = \operatorname{Aut}(L/F)$.
		\item Каждой подгруппе $H$ ставится в соответствие неподвижное подполе $L^H = \{ a \in L \mid \forall h \in H : h(a) = a \}$.
	\end{itemize}
\end{Statement}

\begin{Remark}{}{}
	В общем случае это утверждение верно для расширений Галуа, то есть при условии выполнения равенства $|\operatorname{Aut}(L/K)| = [L : K]$.
\end{Remark}
